% =====================================================================
%  30.4  Linearity of expectation on a rooted tree.
%  Exports: sec:c32-expect  def:c32-ind  lem:c32-indexp  thm:c32-sum
%  prob:c32-game  fig:c32-game  lst:c32-game  lst:c32-depth
%  Solved CF 280C "Game on Tree" (2200, L4), live-verified against the CF
%  problem page (title "C. Game on Tree", 1s TL, 256 MB, n<=1e5, answer
%  tolerance 1e-6).  Two official samples: n=2 edge 1-2 -> 1.5; n=3 star
%  1-2, 1-3 -> 2.0.  Formula E = sum_v 1/depth(v), depth counted in
%  vertices on the root-to-v path so depth(root)=1.
% =====================================================================
\section{Linearity of Expectation on a Rooted Tree}\label{sec:c32-expect}

\begin{previouslybox}[Counting becomes expectation with one small twist]
\(\secref{c32-prufer}\) through \(\secref{c32-matrix}\) produce an integer
answer (often modulo a prime).  This section produces a real: an expectation.
The trick is the same trick a probabilist uses everywhere -- split a random
variable into indicators, and swap the sum with the expectation.  On trees
the split is naturally \emph{per vertex}, and the resulting probability is a
reciprocal of a depth, so \(\secref{c32-expect}\) hands you an \(O(n)\) DFS
where a naive recursion on the whole process would be exponential.
\end{previouslybox}

\begin{topicintro}
Let \(T\) be a rooted tree with root \(r=1\).  A random process on \(T\) has
a state space that grows exponentially with \(n\); writing down a recursion
for \(\mathbb{E}[\text{steps}]\) is hopeless.  But expectation is \emph{linear
without independence} (\(\lemref{c32-indexp}\)): for any random variables
\(X_1,\dots,X_n\) (dependent or not), \(\mathbb{E}[\sum X_i]=\sum\mathbb{E}[X_i]\).
Decompose "number of steps" as \(\sum_v X_v\), where \(X_v=1\) iff \(v\) is
\emph{directly picked} by the process, and each \(\mathbb{E}[X_v] =
\Pr[X_v=1]\) becomes a small rational (\(\thmref{c32-sum}\)).  On the CF 280C
process -- repeatedly choose a remaining vertex uniformly and delete its
subtree -- the answer collapses to \(\sum_v 1/\mathrm{depth}(v)\), a single
DFS.

CF 280C (\(\probref{c32-game}\)) is exactly that computation.  It is a
\(\lv{4}\) problem not because the mathematics is hard but because the
\emph{decomposition} is not obvious -- most first attempts try to write a
DP over the process, which is exponential.

\textbf{Prerequisites.} Ch.~12 (\chref{c09}) for the depth array; basic
probability (uniform distribution over a finite set).  Nothing else.

\textbf{After this section you can}
\begin{itemize}[nosep]
  \item state linearity of expectation without an independence hypothesis; \hfill\lv{3}
  \item decompose a step-count into per-vertex indicators; \hfill\lv{4}
  \item compute \(\Pr[v\text{ is picked}]\) via a symmetry argument; \hfill\lv{4}
  \item solve CF 280C in one DFS. \hfill\lv{4}
\end{itemize}
\end{topicintro}

% ---------------------------------------------------------------
\subsection{Linearity without independence}
% ---------------------------------------------------------------

\begin{definition}{Indicator variable of an event}{c32-ind}
For an event \(A\) in a probability space \(\Omega\), its \term{indicator}
\(\mathbf{1}_A:\Omega\to\{0,1\}\) is \(\mathbf{1}_A(\omega)=1\) iff
\(A\) occurs at \(\omega\).  \(\mathbb{E}[\mathbf{1}_A]=\Pr[A]\).
\end{definition}

\begin{lemma}{Linearity of expectation}{c32-indexp}
For random variables \(X_1,\dots,X_n\) on a common probability space, not
assumed independent, \(\mathbb{E}\!\left[\sum_i X_i\right]=\sum_i
\mathbb{E}[X_i]\).
\end{lemma}

\begin{proofbox}
Follows from \(\mathbb{E}[X]=\sum_{\omega} P(\omega) X(\omega)\) and the
distributive law: swap the order of the two finite sums.  No independence
assumption is used. \(\Box\)
\end{proofbox}

\begin{theorem}{Depth-reciprocal sum for CF 280C}{c32-sum}
Let \(T\) be a rooted tree on \([n]\) with root \(1\).  Consider the process
that repeatedly picks a uniformly random remaining vertex and deletes its
subtree, until the tree is empty.  The expected number of steps is
\[
  \mathbb{E}[\text{steps}] \;=\; \sum_{v=1}^{n} \frac{1}{\mathrm{depth}(v)},
\]
where \(\mathrm{depth}(v)\) is the number of vertices on the root-\(v\)
path, so \(\mathrm{depth}(1)=1\).
\end{theorem}

\begin{proofbox}
The process is equivalent to assigning each vertex \(v\) an independent
uniform random priority \(\pi_v\in[0,1]\), then processing vertices in
increasing priority order and marking a step for every vertex whose
ancestors on the root-path have all been processed (i.e., were the "next
chosen" at their own moment).  A vertex \(v\) contributes one step iff
\(\pi_v\) is the minimum on the root-\(v\) path; by symmetry over the
\(\mathrm{depth}(v)\) path vertices, \(\Pr[\pi_v\text{ is the min}]=1/\mathrm{depth}(v)\).

Now write \(X_v = \mathbf{1}\{v\text{ contributes a step}\}\).  Total steps
\(N = \sum_v X_v\).  By \(\lemref{c32-indexp}\),
\(\mathbb{E}[N] = \sum_v \Pr[X_v=1] = \sum_v 1/\mathrm{depth}(v)\). \(\Box\)
\end{proofbox}

\begin{lstlisting}[caption={Root a tree at 1 and compute depth via a single DFS or BFS; the expectation is the sum of reciprocal depths. Iterative DFS with an explicit stack avoids recursion-depth issues on n up to 1e5. ASCII only, identifiers without underscores}, label={lst:c32-depth}]
#include <bits/stdc++.h>
using namespace std;
int main(){
  int n; if(scanf("%d",&n)!=1) return 0;
  vector<vector<int> > adj(n+1);
  for(int i=0;i<n-1;i++){
    int u,v; scanf("%d %d",&u,&v);
    adj[u].push_back(v); adj[v].push_back(u);
  }
  vector<int> depth(n+1,0);
  depth[1]=1;
  vector<int> st; st.push_back(1);
  vector<int> par(n+1,0);
  while(!st.empty()){
    int u=st.back(); st.pop_back();
    for(int v : adj[u]) if(v!=par[u]){
      par[v]=u; depth[v]=depth[u]+1; st.push_back(v);
    }
  }
  double E=0.0;
  for(int v=1; v<=n; v++) E += 1.0/depth[v];
  printf("%.20f\n", E);
  return 0;
}
\end{lstlisting}

% ---------------------------------------------------------------
\subsection{Solved: CF 280C -- Game on Tree (2200, L4)}
% ---------------------------------------------------------------

\solvedproblem{c32-game}{Game on Tree}{CF 280C}{https://codeforces.com/problemset/problem/280/C}{CF 2200 L4}{Linearity of expectation (\(\lemref{c32-indexp}\)); per-vertex indicator decomposition; one DFS for depth (\(\thmref{c32-sum}\), \(\coderef{c32-depth}\))}

\begin{problemstatement}[CF 280C, verified against the live CF problem page]
A rooted tree on \(n\) vertices has root \(1\); the next \(n-1\) lines
describe edges as pairs \((a_i,b_i)\) with \(1\le a_i,b_i\le n\) and
\(a_i\ne b_i\), and the graph is guaranteed to be a tree.  A one-player game
runs as follows: while at least one vertex remains, pick a uniformly random
vertex among those remaining and delete its whole rooted subtree (the vertex
and every descendant of it in the residual forest).  The game ends when no
vertices remain.  Report the expected number of steps as a real number to
absolute or relative error \(\le 10^{-6}\).  Constraint: \(1\le n\le
10^5\).  Samples: input \(2,1\;2\) outputs \(1.5\); input \(3,1\;2,1\;3\)
outputs \(2.0\).
\end{problemstatement}

\begin{figure}[H]
\centering
\begin{tikzpicture}[diagram, scale=0.9, transform shape,
   tnode/.style={circle, draw=primaryblue!70, fill=primaryblue!8, line width=0.7pt,
                 minimum size=6mm, font=\scriptsize\bfseries, inner sep=1pt}]
  \node[tnode] (r) at (0,1.2) {1};
  \node[tnode] (a) at (-1.2,-0.2) {2};
  \node[tnode] (b) at (1.2,-0.2) {3};
  \draw[primaryblue!70, line width=1pt] (r)--(a) (r)--(b);
  \node[annot, font=\scriptsize, anchor=north] at (0,-1.4)
    {sample 2: depths \((1,2,2)\), \(E = 1 + 1/2 + 1/2 = 2\)};
\end{tikzpicture}
\caption{Sample 2 of CF 280C.  Root \(1\) has depth \(1\); vertices \(2\) and
\(3\) have depth \(2\).  The depth-reciprocal sum gives \(E=2\), matching the
published sample output \(2.00000000000000000000\).}
\label{fig:c32-game}
\end{figure}

\begin{attemptgate}[Attempt it yourself first -- \lv{4}, give yourself 25 minutes]
Before coding: can you write a recursion on "expected steps from a residual
forest \(F\)"?  What goes wrong, and what decomposition avoids it?
\peek{A recursion on \(F\) has \(2^n\) states; memoising is not an option.
Instead, ask: what is one \emph{step} made of?  It picks exactly one vertex,
so the total number of steps equals the number of vertices DIRECTLY picked
(never covered by an ancestor).}
\nudge{Call that event \(A_v\).  The total steps \(N = \sum_v
\mathbf{1}_{A_v}\).  Linearity of expectation (\(\lemref{c32-indexp}\)) gives
\(\mathbb{E}[N] = \sum_v \Pr[A_v]\), and the sum runs over \(n\) terms with
no independence required.}
\fullhint{The process on \(T\) is equivalent to assigning each vertex an
i.i.d.\ continuous uniform priority and processing vertices in increasing
priority order; a vertex \(v\) is "directly picked" iff its priority is the
minimum on the root-\(v\) path.  By symmetry over the path's
\(\mathrm{depth}(v)\) labels, this has probability \(1/\mathrm{depth}(v)\).
Compute depths with one DFS (\(\coderef{c32-depth}\)) and sum the
reciprocals; \(\bigO{n}\) time, \(\bigO{n}\) space.}
\end{attemptgate}

\solstep{Understanding}
CF 280C looks like a dynamic-programming-on-subtrees problem because the
input is a rooted tree, but the state is a \emph{set of remaining vertices},
which explodes exponentially.  The rescue is \(\lemref{c32-indexp}\): you
never need to reason about the whole process, only about whether each single
vertex is picked directly.  The indicator decomposition collapses the
exponential state space to \(n\) independent questions, each of which has
the same answer for the same-depth vertices.  What remains is a formula,
and the formula is one DFS.

The priority-ordering equivalence spelled out in the full hint above is standard: drawing
each vertex uniformly at random from \([0,1]\) and taking them in increasing
order is a uniformly random permutation of the vertices.  The process
"pick a uniformly random remaining vertex at each step" is a uniformly
random permutation as well, and it deletes a vertex's subtree at the
vertex's own moment in the permutation.

\solstep{Key Observations}
\begin{observation}
Every step picks exactly one vertex that survives until that step, and
deletes a set of vertices (its subtree).  The vertex \emph{doing} the
deletion at step \(t\) is called the "picked" vertex; every other deleted
vertex is "covered" by an ancestor's pick.  Total steps \(=\) total picked
vertices \(=\sum_v \mathbf{1}_{A_v}\).
\end{observation}
\begin{observation}
\(A_v\) holds iff \(v\) appears before any of its ancestors in the random
permutation.  Equivalently, among the \(\mathrm{depth}(v)\) vertices on the
root-to-\(v\) path, \(v\) is the earliest.  By symmetry of a uniformly
random permutation, each of these \(\mathrm{depth}(v)\) labels is equally
likely to be the earliest, so \(\Pr[A_v]=1/\mathrm{depth}(v)\).
\end{observation}
\begin{observation}
No independence between \(A_u\) and \(A_v\) is required.  Linearity of
expectation (\(\lemref{c32-indexp}\)) sums \(n\) rationals regardless of the
dependence structure among the indicators.
\end{observation}

\solstep{Think It With Me}
\begin{thinkbox}
\deadend{Recursion on subtrees: \(\mathbb{E}[T] = 1 + \frac{1}{|V(T)|}
\sum_v \mathbb{E}[T\setminus \mathrm{subtree}(v)]\)}{The state space is the
set of residual forests reachable by subtree deletions; there are
exponentially many, and the recursion has no memoisable subproblem.}
\deadend{Try a symmetry argument over "expected depth"}{That measures only a single vertex's path contribution, not the sum of step indicators, so it undercounts the total expectation.}
\deadend{Write a DP with state "root picked or not"}{That flag captures whether the root has fired, but the residual forest below the root still has exponentially many configurations once you recurse.}
\aha{Split into indicators, one per vertex.  Linearity applies with no
independence, and each indicator's probability is \(1/\mathrm{depth}(v)\)
by symmetry of a uniform permutation.}
\end{thinkbox}

\solstep{The Approach Ladder}
\begin{center}\renewcommand{\arraystretch}{1.25}\setlength{\tabcolsep}{3pt}
\begin{tabular}{@{}c p{0.26\linewidth} p{0.30\linewidth} l c@{}}
\toprule
\textbf{\#} & \textbf{Approach} & \textbf{Idea} & \textbf{Time} & \textbf{Verdict}\\
\midrule
1 & \lv{2}\ full recursion & exponential state space & \(\bigO{2^n}\) & \tle\\
2 & \lv{3}\ DP with "root picked" flag & still exponential residual forest & \(\bigO{2^n}\) & \tle\\
3 & \lv{4}\ \textbf{indicator + linearity} & per-vertex symmetry & \(\bigO{n}\) & \ac\\
\bottomrule
\end{tabular}\end{center}
Rungs 1 and 2 are natural but doomed.  Rung 3 is one DFS and one sum.

\solstep{Solution}
Root the tree at vertex \(1\) (which the statement already fixes).  Run one
DFS to fill \(\mathrm{depth}[v]\) for every \(v\) (\(\coderef{c32-depth}\));
the recursion is \(O(n)\) total.  Then compute \(\sum_v 1/\mathrm{depth}[v]\)
as a \code{double} and print with sufficient precision (\texttt{\%.20f} is a
safe default for a \(10^{-6}\) error budget).

\solstep{Implementation}
\begin{lstlisting}[caption={CF 280C. Iterative DFS for depth, sum reciprocal depths, print with wide precision. The iterative stack avoids recursion-depth issues at n=1e5. ASCII only, identifiers without underscores}, label={lst:c32-game}]
#include <bits/stdc++.h>
using namespace std;
int main(){
  int n; if(scanf("%d",&n)!=1) return 0;
  vector<vector<int> > adj(n+1);
  for(int i=0;i<n-1;i++){
    int u,v; scanf("%d %d",&u,&v);
    adj[u].push_back(v); adj[v].push_back(u);
  }
  vector<int> depth(n+1,0), par(n+1,0);
  vector<int> st; st.push_back(1); depth[1]=1;
  while(!st.empty()){
    int u=st.back(); st.pop_back();
    for(int v : adj[u]) if(v!=par[u]){
      par[v]=u; depth[v]=depth[u]+1; st.push_back(v);
    }
  }
  double E=0.0;
  for(int v=1; v<=n; v++) E += 1.0/depth[v];
  printf("%.20f\n", E);
  return 0;
}
\end{lstlisting}

\solstep{Proof of Correctness}
\begin{proofbox}
\textbf{(1)} The process is equivalent to a uniformly random permutation of
\([n]\): at each step pick the surviving vertex with the smallest position
in the permutation, and delete its subtree in the surviving forest.  This
generates the same distribution over "which vertex is picked at each step"
as the original process, because at every step the choice among the
survivors is uniform by the permutation's symmetry.

\textbf{(2)} A vertex \(v\) is picked by this permutation-protocol iff \(v\)
has the smallest permutation position among all vertices on the root-\(v\)
path in the original tree.  Reason: if an ancestor \(a\) of \(v\) has a
smaller position, then \(a\) is picked first and its subtree-deletion
removes \(v\), so \(v\) is never directly picked.  Conversely if \(v\) has
the smallest position on the path, then no ancestor of \(v\) has fired
before \(v\) does, so \(v\) survives until its own turn.

\textbf{(3)} By symmetry of a uniformly random permutation restricted to
any fixed set of \(\mathrm{depth}(v)\) positions, the smallest of them is
uniformly distributed over those \(\mathrm{depth}(v)\) vertices, so
\(\Pr[A_v] = 1/\mathrm{depth}(v)\).

\textbf{(4)} \(N=\sum_v \mathbf{1}_{A_v}\), so by \(\lemref{c32-indexp}\)
(linearity without independence) \(\mathbb{E}[N]=\sum_v 1/\mathrm{depth}(v)\),
which is what \(\coderef{c32-game}\) computes. \(\Box\)
\end{proofbox}

\solstep{Trace on the Official Samples}
\begin{tracebox}[Sample 1: \, \(n=2\), edge \(1\!-\!2\)]
Depths \(\mathrm{depth}[1]=1\), \(\mathrm{depth}[2]=2\).  Sum \(=1+1/2=
1.5\).  The published sample output is \(1.50000000000000000000\).
Direct check: two runs of the game -- either pick \(1\) first (probability
\(1/2\), \(1\) step; the tree empties immediately) or pick \(2\) first
(probability \(1/2\), \(1\) step deletes \(2\), then one more step deletes
residual \(1\), so \(2\) steps).  \(\mathbb{E}=1/2\cdot 1 + 1/2\cdot 2 =
1.5\).  Matches.
\end{tracebox}
\begin{tracebox}[Sample 2: \, \(n=3\), edges \(1\!-\!2\) and \(1\!-\!3\)]
Depths \((1,2,2)\).  Sum \(=1+1/2+1/2=2\).  Published sample output
\(2.00000000000000000000\).  Direct: pick root \(1\) (probability \(1/3\))
ends in \(1\) step; else pick \(2\) or \(3\) first (probability \(2/3\)),
which deletes one leaf; the residual tree is an edge on \(\{1,3\}\) or
\(\{1,2\}\), expected \(1.5\) more steps; total \(1/3\cdot 1 + 2/3\cdot
(1 + 1.5) = 1/3 + 5/3 = 2\).  Matches.
\end{tracebox}

\solstep{Complexity}
\begin{complexitybox}
Building the adjacency list: \(\bigO{n}\).  One iterative DFS fills
\(\mathrm{depth}[\cdot]\): \(\bigO{n}\).  Sum of reciprocals: \(\bigO{n}\).
Total \(\bigO{n}\) time, \(\bigO{n}\) space.  For \(n=10^5\) this is
comfortably under the 1-second limit.  The \code{double} accumulator has a
round-off error of \(\bigO{n\epsilon_{\mathrm{machine}}}\approx 10^{-11}\)
for \(n=10^5\), far inside the \(10^{-6}\) tolerance.
\end{complexitybox}

\begin{hardeningbox}
\begin{itemize}[nosep,leftmargin=1.4em]
\item \textbf{Iterative DFS.}  \(n\le 10^5\), and a chain of that length
      recurses \(10^5\) frames deep -- an easy stack overflow with default
      settings.  The explicit-stack form in \(\coderef{c32-game}\) is
      safer and equally short.
\item \textbf{Depth starts at 1.}  A vertex at BFS level 0 (the root) is at
      path-length \(\mathrm{depth}=1\) (root itself).  Off-by-one here
      changes the answer by exactly \(1\) on the sample \(n=2\) test.
\item \textbf{Print format.}  \texttt{\%.20f} is overkill but safe; do not
      switch to \texttt{\%g}, which may print scientific notation and can
      drop significant digits under some libraries.
\item \textbf{Adjacency vector of vectors.}  Fine for \(n\le 10^5\); for
      tighter memory budgets use a flat edge list with head/next arrays.
\end{itemize}
\end{hardeningbox}

\begin{recallbox}[30.4 in one breath]
\emph{Every step picks exactly one vertex; total steps is \(\sum_v
\mathbf{1}_{A_v}\).  Linearity of expectation (\(\lemref{c32-indexp}\))
sums the \(\Pr[A_v]\) with no independence assumption.  Symmetry of a
uniform permutation on a fixed path gives \(\Pr[A_v]=1/\mathrm{depth}(v)\).
One DFS, one sum, done (\(\thmref{c32-sum}\), \(\coderef{c32-game}\)).}
\selftest{
  \item State linearity of expectation and name one hypothesis it does
        \emph{not} require.
  \item Explain why the "random permutation" reformulation of CF 280C
        preserves the process's distribution.
  \item Compute the answer by hand for a path \(1\!-\!2\!-\!3\) rooted at \(1\).
        \emph{(Depths \((1,2,3)\), sum \(1+1/2+1/3 = 11/6\).)}
}
\end{recallbox}
